\mychapter{Hardware}{Hardware}{}
\label{chap:hardware}

\tikzset{
  blk/.style={draw, rounded corners=2pt, align=center, font=\footnotesize,
              minimum width=2.6cm, minimum height=0.95cm, inner sep=3pt},
  mdl/.style={blk, fill=black!8},
  hw/.style={blk, dashed},
  st/.style={draw, rounded corners=2pt, align=center, font=\footnotesize,
             minimum width=2.7cm, minimum height=0.9cm},
  arr/.style={-{Stealth[length=2mm]}, thick},
  lbl/.style={font=\scriptsize, fill=white, inner sep=1pt}
}

In the design Chapter~\ref{chap:design} was where the conversation went about how to split
the decisions between the rover and the brain. In this chapter we will discuss how the rover
connects to the brain and how it is controlled by the brain. Translating the plan from the planner
into actions a rover can execute. The rover is an off-the-shelf Hiwonder advanced turboPi kit. It uses the
open source Robot operating system 2 to execute motion commands. 


\mysection{The Turbopi platform}{The TurboPi Platform}
\label{sec:turbopi_platform}

The rover is a Hiwonder TurboPi, Advanced kit, built around a Raspberry
Pi~5. The only actuators on this specific unit are a four-wheel mecanum
base, which can strafe sideways as well as drive and turn, and a
two-servo pan/tilt gimbal carrying the camera. A ultrasonic sensor (sonar) with an RGB ring, 
sits on the front of the chassis, and a row of infrared
line-detection sensors sits underneath the chassis. There is no arm and no
gripper on this build.

The Advanced kit does not run a lightweight Python SDK, which is usually on the base robot kit.
It can call motor functions directly. Instead it runs ROS2 inside a Docker
container on the Pi, and every piece of hardware, driving, the gimbal,
the sonar RGB ring, is reached through a ROS2 topic or service rather
than a function call. As the rover starts up it lauches \texttt{bringup.launch.py},
It launches the \texttt{mecanum} node which can turn the velocity command into wheel speeds. 
Then the important \texttt{ros\_robot\_controller} which. talks to the motor and servo hardware. 
There is a serial link connected to a STM32 co-processor on the expansion board, a camera connected over USB,
a small HTTP video server and  The main contact point into the rover is over the \texttt{rosbridge\_websocket}.


\begin{figure}[htbp]
  \centering
  \begin{tikzpicture}[node distance=1.1cm]
    \node[blk] (brain) at (0,0) {Brain\\(laptop)};
    \node[hw] (bridge) at (5.2,1.1) {\texttt{rosbridge\_}\\\texttt{websocket}\\:9090};
    \node[hw] (video) at (5.2,-1.1) {\texttt{web\_video\_}\\\texttt{server}\\:8080};
    \node[mdl] (mec) at (9.6,1.9) {\texttt{mecanum}};
    \node[mdl] (cam) at (9.6,0.3) {\texttt{usb\_cam}};
    \node[mdl] (rrc) at (9.6,-1.5) {\texttt{ros\_robot\_}\\\texttt{controller}};
    \node[hw] (stm) at (13.4,-1.5) {STM32\\(serial)};
    \draw[arr] (brain.east) -- node[lbl,above]{roslibpy} (bridge.west);
    \draw[arr] (brain.east) -- node[lbl,below]{HTTP} (video.west);
    \draw[arr] (bridge) -- (mec);
    \draw[arr] (video) -- (cam);
    \draw[arr] (bridge.east) -- (rrc.west);
    \draw[arr] (mec) -- (rrc);
    \draw[arr] (rrc) -- (stm);
  \end{tikzpicture}
  \caption{What the brain actually reaches on the Pi. Everything above
  \texttt{rosbridge\_websocket} and \texttt{web\_video\_server} is
  Hiwonder's stock ROS2 stack, started as one unit by
  \texttt{bringup.launch.py} inside the Pi's Docker container.}
  \label{fig:hw-bridge}
\end{figure}


\mysection{The bridge the brain actually uses}{The Bridge the Brain Actually Uses}
\label{sec:hw_bridge}


The \texttt{rosbridge\_websocket} is a stock ROS2 package that exposes topics
and services as plain JSON over a WebSocket. The brain never needs to read DDS stack or
and of Hiwonders Actual packages for control. This was chosen over the brain running \texttt{rclpy}
since the brain already has existing WebSockets for \texttt{/ws/dialogue} and \texttt{/ws/execution}

To actually drive the rover you need to send values to the \texttt{geometry\_msgs/Twist} message on the
\texttt{cmd\_vel} topic. \texttt{linear.x} and \texttt{linear.y} are
forward and strafe speed, and \texttt{angular.z} is turn rate. The angular term is summed into each of the
wheels kinematics before scaling to a motor speed, so it is not treat as a radian per second value. To figure out
what range of values actually turn the rover accurately I chose a range of 3 to 8 which is used in the Hiwonder's
own stock code. The gimbal is driven the same way, a
\texttt{SetPWMServoState} message on
\texttt{/ros\_robot\_controller/pwm\_servo/set\_state} naming a servo id,
1 for tilt and 2 for pan on this unit, and a target position around a
centre of 1500.

The frames from the camera are pulled directly from the websocket running on the Pi which had a port of 8080. The
endpoint was at \texttt{web\_video\_server}.


%ask lecturere if I need to add this section
\mysection{Reusing the stock line follower}{Reusing the Stock Line Follower}
\label{sec:hw_lf}
